\chapter{Conceptos comunes abordados en grado y TFG}
\label{cap:capitulo12}

\vspace{0.25cm}

\section{Introducción}
\label{sec:segundaseccion}

A lo largo del Grado en Ingeniería de Robótica Software, se ha adquirido una amplia variedad de conocimientos en múltiples disciplinas, desde fundamentos matemáticos y físicos hasta áreas como la programación, la robótica, los Sistemas Distribuidos o la Inteligencia Artificial.

\vspace{0.25cm}

Es por ello que este TFG no se encuentra aislado de lo visto en el grado, ya que integra gran parte de todos estos conocimientos, aplicándolos de forma conjunta para hacer posible el diseño, implementación y validación de un sistema robótico real. Dicho de esta manera, se puede entender que este TFG actúa como punto de unión con gran parte de los conceptos aprendidos durante toda la formación académica.

\vspace{0.25cm}

En este capítulo se explican todas las relaciones entre las asignaturas del Grado en Ingeniería de Robótica Software y los distintos componentes de este TFG, destacando cómo y en qué medida ha influido cada una de ellas en el resultado final.

\section{Relación con asignaturas de primer curso}
\label{sec:segundaseccion}

Las asignaturas de primer curso han proporcionado los fundamentos necesarios sobre los que se apoya todo el desarrollo posterior de este TFG.

\vspace{0.25cm}

Las asignaturas de \textit{Álgebra}, \textit{Cálculo} y \textit{Fundamentos Físicos de la Robótica}  han resultado esenciales para poder formular el modelo cinemático del sistema. El uso de vectores, distancias euclídeas, relaciones geométricas, magnitudes físicas de posición y movimiento, y relaciones espaciales, han permitido describir la posición del efector final en el espacio y establecer relaciones entre todas las variables físicas que componen el sistema.

\vspace{0.25cm}

Y por otro lado, asignaturas como \textit{Fundamentos de la Programación} y \newline \textit{Algoritmos y Estructuras de Datos} han sido claves para llevar a cabo el desarrollo software del sistema, donde la capacidad de estructurar código, definir funciones y diseñar algoritmos, ha resultado útil para implementar el modelo cinemático y la planificación de trayectorias.

\section{Relación con asignaturas de segundo curso}
\label{sec:segundaseccion}

Las asignaturas de segundo curso introducen conceptos orientados al diseño de sistemas y al ámbito específico de la robótica, las cuales establecen una relación directa en el desarrollo de este TFG.

\vspace{0.25cm}

La asignatura de \textit{Diseño Software} ha influido directamente en la organización modular del sistema desarrollado, donde la estructuración del código y la claridad en la arquitectura del sistema reflejan los conceptos vistos en esta asignatura.

\vspace{0.25cm}

Por otro lado, la asignatura de \textit{Sensores y Actuadores} se relaciona directamente con el desarrollo del sistema hardware, concretamente en el uso de los servomotores encargados de accionar las poleas, permitiendo comprender cómo se traduce una señal de control en un movimiento físico real.

\vspace{0.25cm}

Y por último, las asignaturas de \textit{Arquitecturas Software para Robots} y \newline \textit{Sistemas Operativos} se relacionan directamente con la integración del sistema en ROS 2, ya que introducen conceptos como la comunicación entre componentes, el diseño de sistemas robóticos complejos, la gestión de recursos y la ejecución concurrente.

\section{Relación con asignaturas de tercer curso}
\label{sec:segundaseccion}

Las asignaturas de tercer curso son las que constituyen el núcleo de este TFG, especialmente en lo que respecta a robótica y sistemas software avanzados.

\vspace{0.25cm}

Las asignaturas de \textit{Robótica Móvil}, \textit{Modelado y Simulación de Robots}, y \newline \textit{Sistemas Empotrados y de Tiempo Real} tienen una relación directa con el desarrollo del modelo cinemático y su validación, permitiendo comprender cómo modelar el movimiento de un robot y cómo simular su comportamiento para posteriormente integrarlo en un sistema real, donde aspectos como la ejecución en tiempo real y la interacción con dispositivos físicos cobran una gran importancia.

\newpage

Por otro lado, la asignatura de \textit{Sistemas Distribuidos y Concurrentes} refleja claramente el uso de ROS 2, concretamente en la organización del sistema en nodos que se comunican entre sí a través de topics, todo ello mientras se sigue un modelo distribuido.

\vspace{0.25cm}

Y por último, la asignatura de \textit{Planificación y Sistemas Cognitivos} posee una relación directa con la generación de trayectorias, ya que permite al sistema generar movimientos contínuos a partir de múltiples posiciones objetivo.

\section{Relación con asignaturas de cuarto curso}
\label{sec:segundaseccion}

Las asignaturas de cuarto curso aportan una visión más aplicada y especializada, aunque también se pueden ver reflejadas en el desarrollo de este TFG.

\vspace{0.25cm}

La asignatura de \textit{Robótica de Servicio} permite contextualizar el sistema desarrollado en aplicaciones reales, por ejemplo, en sistemas de posicionamiento o de manipulación ligera, donde los CDPR pueden tener un papel relevante.

\vspace{0.25cm}

Y por otro lado, la asignatura de \textit{Mecatrónica} resulta clave en la integración entre hardware y software, uno de los aspectos fundamentales de este TFG, ya que permite plantear un sistema físico controlado a partir del sistema software desarrollado.

\section{Integración global de conocimientos en el TFG}
\label{sec:segundaseccion}

El desarrollo de este TFG es un ejemplo claro de la necesidad de combinar elementos de diferentes áreas para construir un sistema completo y funcional.

\vspace{0.25cm}

Por ello, se han aplicado conceptos matemáticos para definir el modelo cinemático, conocimientos de programación para implementar el modelo diseñado en distintos lenguajes y entornos, conceptos de arquitectura software para integrar el sistema en ROS 2, y conceptos de hardware para implementar el sistema en físico.

\vspace{0.25cm}

Es por todo esto que toda esta integración de conocimientos permite abordar el problema de forma global y comprender cómo interactúan entre sí todos los componentes que forman el sistema completo.

\newpage

\section{Conclusión}
\label{sec:segundaseccion}

Este TFG representa una combinación de gran parte de los conocimientos adquiridos a lo largo del Grado en Ingeniería de Robótica Software.

\vspace{0.25cm}

A través del desarrollo de un sistema basado en un CDPR, se han podido poner en práctica conceptos de múltiples asignaturas, lo que ha permitido demostrar la capacidad de integrar conocimientos teóricos y aplicarlos en un contexto práctico y real.

\vspace{0.25cm}

Es por ello que este proceso no sólo ha permitido consolidar los aprendizajes adquiridos a lo largo del grado, también han permitido desarrollar una visión más completa del ámbito de la robótica, la cual requiere la combinación de diferentes áreas de conocimiento para la resolución de problemas complejos.
